\documentclass[10pt,twocolumn]{article}
\usepackage[utf8]{inputenc}
\usepackage[T1]{fontenc}
\usepackage{amsmath,amssymb}
\usepackage{graphicx}
\usepackage{booktabs}
\usepackage{hyperref}
\usepackage[margin=2cm]{geometry}
\usepackage{caption}
\usepackage{subcaption}
\usepackage{xcolor}
\usepackage{algorithm}
\usepackage{algorithmic}

\newcommand{\pending}[1]{\textcolor{red}{\textbf{[PENDING: #1]}}}

\title{\textbf{TinyDamage: Multi-Scale Contrastive Segmentation for Fine-Grained Vehicle Damage Detection}}

\author{
  Anjali Pai \quad Nitya Ravi \quad Vishwajeet Shivaji Hogale \\
  Cornell University \\
  \texttt{\{ap, nr, hogale.v\}@cornell.edu}
}

\date{}

\begin{document}
\maketitle

% ============================================================
\begin{abstract}
Segmenting tiny vehicle damages like scratches, hairline cracks, and minor dents, is critical for automated insurance claim processing yet remains a persistent challenge. These defects occupy few pixels, produce weak gradient signals during training, and are easily confused with surface reflections and paint textures. While recent work has advanced car damage segmentation through increasingly sophisticated architectures, we show that the choice of loss function and optimization strategy has an outsized and underexplored effect on tiny-damage detection-and that naively applying popular losses designed for class imbalance can catastrophically suppress small-object predictions.

We propose \textbf{TinyDamage}, a multi-scale segmentation framework that integrates three components designed to address the specific failure modes of tiny-damage learning: (1)~a Feature Pyramid Network (FPN) encoder for multi-scale feature preservation, (2)~a Tiny-Object Contrastive Module that learns to separate damage embeddings from visually similar background features in a projected pixel embedding space, and (3)~a Gradient-Aware Boundary Module that provides controlled edge supervision without dominating the training signal. Our framework is architecture-agnostic and evaluated across three segmentation backbones: U-Net, YOLOv8-Seg, and Mask2Former.

Through systematic ablation on the CarDD dataset, we first establish that focal loss and gradient-dominant objectives-despite their popularity-collapse tiny-damage detection to zero when applied naively, while a simple CE+Dice baseline with cosine annealing provides a surprisingly strong foundation (test $\text{DET}_l = 1.000$, test mIoU $= 0.613$). We then show that TinyDamage's contrastive and gradient-aware modules, integrated as controlled auxiliary objectives, recover boundary precision without sacrificing detection recall. We introduce $\text{DET}_l$, a per-category metric for evaluating tiny-damage detection independently of aggregate overlap, and argue for its adoption alongside standard mIoU in this domain.
\end{abstract}

% ============================================================
\section{Introduction}
\label{sec:intro}

Vehicle damage assessment is a bottleneck in the insurance claims pipeline. Manual inspection by trained adjusters is costly, subjective, and slow-yet the industry processes millions of claims annually. Automating even the initial triage of damage documentation could substantially reduce processing time and improve consistency. Computer vision methods, particularly semantic and instance segmentation, offer a viable path toward this goal.

Recent years have seen significant architectural progress in vehicle damage segmentation. Purpose-built datasets such as CarDD~\cite{wang2023cardd} provide high-quality annotations for benchmarking. Specialized architectures including MARS~\cite{panboonyuen2023mars}, which refines instance masks through quadtree-based self-attention, ALBERT~\cite{panboonyuen2025albert}, which jointly segments damage types and vehicle parts using a multi-headed transformer, and SLICK~\cite{panboonyuen2025slick}, which distills these capabilities into a real-time model, have pushed the state of the art on standard metrics. Foundation model adaptations such as PerSAM~\cite{zhang2023persam} have also been explored for one-shot personalized segmentation tasks that could extend to damage inspection.

However, a fundamental challenge persists across these advances: \textbf{tiny damages are hard.} Scratches, hairline cracks, and shallow dents occupy a small fraction of image pixels, produce weak supervisory signal during training, and share visual characteristics with benign surface features-reflections, paint texture gradients, and lighting artifacts. Standard evaluation metrics like mIoU and mAP mask this problem because they are dominated by large, visually salient categories such as crushed panels and shattered glass.

The segmentation community has addressed small-object challenges primarily through architectural innovation: multi-scale feature pyramids, attention mechanisms, and specialized detection heads. A complementary but underexplored direction is the \textbf{optimization landscape}-the loss functions, schedulers, and training dynamics that determine whether a model learns to attend to tiny foreground regions at all. This direction is especially relevant because loss and scheduler choices are the most accessible levers for practitioners.

In this work, we make two contributions. First, through a systematic baseline ablation, we reveal that popular optimization choices can \textbf{catastrophically fail} for tiny-damage segmentation: focal loss, despite being designed for class imbalance, collapses tiny-damage detection to zero; gradient-based boundary supervision similarly collapses when used as a dominant objective. These findings motivate the design of our second contribution: \textbf{TinyDamage}, a multi-scale segmentation framework whose auxiliary modules are specifically designed to avoid these failure modes while improving feature discriminability and boundary precision for small damages.

Our contributions are:
\begin{enumerate}
    \item A \textbf{baseline ablation study} revealing that focal-dominant and gradient-dominant loss configurations collapse tiny-damage detection, while CE+Dice with cosine annealing provides a strong foundation (Section~\ref{sec:ablation}).
    \item The \textbf{TinyDamage framework}, integrating FPN-based multi-scale encoding, a Tiny-Object Contrastive Module, and a Gradient-Aware Boundary Module as controlled auxiliary objectives (Section~\ref{sec:method}).
    \item \textbf{Cross-architecture evaluation} on U-Net, YOLOv8-Seg, and Mask2Former, demonstrating consistent improvements across architectural paradigms (Section~\ref{sec:cross_arch}).
    \item $\textbf{DET}_l$, a per-category detection metric for tiny-damage evaluation (Section~\ref{sec:metrics}).
\end{enumerate}

% ============================================================
\section{Related Work}
\label{sec:related}

\subsection{Vehicle Damage Segmentation}

Vision-based vehicle damage assessment has progressed from image classification to pixel-level segmentation. Early detection-based methods used Faster R-CNN and YOLO variants for bounding-box localization, while instance segmentation approaches such as Mask R-CNN~\cite{he2017maskrcnn} brought pixel-level accuracy.

The CarDD dataset~\cite{wang2023cardd} provides over 4,000 high-resolution images with 9,000+ annotated instances across six damage categories (scratch, dent, crack, glass shatter, lamp broken, tire flat) and has become a standard benchmark. MARS~\cite{panboonyuen2023mars} introduced self-attention with sequential quadtree nodes for mask refinement, achieving +2.3 maskAP over Mask R-CNN on a Thai car-damage dataset. ALBERT~\cite{panboonyuen2025albert} extends transformer-based segmentation to jointly handle 26 real damage types, 7 fake damage variants, and 61 vehicle parts through multi-headed classification. SLICK~\cite{panboonyuen2025slick} distills ALBERT into a lightweight model with 7$\times$ faster inference via selective localization and knowledge fusion modules.

A recent systematic review~\cite{hasan2025survey} catalogues the progression from classification through detection to segmentation methods and identifies fine-grained damage delineation as a key open challenge. PerSAM~\cite{zhang2023persam} demonstrates training-free personalization of the Segment Anything Model using a single reference image, a paradigm with potential applicability to damage-specific segmentation. Furthermore, Vision-Language Models (VLMs) have emerged as strong contenders due to their advanced reasoning capabilities. In our preliminary experiments with the Qwen-VL 2B model, we found it to be extremely powerful at diagnosing vehicle issues; however, it struggled to produce accurate segmentation masks to localize where the damage was. While fine-tuning VLMs on damage datasets could potentially bridge this localization gap, the heavy computational cost associated with their attention mechanisms makes them less practical than dedicated segmentation architectures for this task.

While these methods advance architectural design, the impact of loss function and optimizer choices on tiny-damage detection has not been systematically studied in this domain.

\subsection{Small-Object Segmentation and Loss Design}

Small-object segmentation is a cross-domain challenge addressed through both architectural and loss-level solutions. Feature Pyramid Networks (FPN)~\cite{lin2017fpn} provide multi-scale representations that preserve fine-grained spatial detail.

At the loss level, focal loss~\cite{lin2017focal} reweights cross-entropy to focus on hard examples via a modulating factor $(1-p_t)^\gamma$. Dice loss~\cite{milletari2016vnet} optimizes region overlap directly and is less sensitive to class imbalance. Boundary-aware losses penalize edge misalignment using gradient-based supervision. Composite losses combining multiple objectives are widely used but their interactions are rarely studied in isolation.

\subsection{Contrastive Learning for Segmentation}

Contrastive learning has been adapted from representation learning to dense prediction tasks. Pixel-level contrastive losses learn embeddings where same-class pixels cluster together and different-class pixels separate, improving feature discriminability for challenging classes. Recent work applies contrastive objectives as auxiliary losses alongside standard segmentation losses, with benefits for rare and small-object categories.

Our Tiny-Object Contrastive Module draws on this paradigm but specifically targets the failure mode we identify in our ablation: the tendency of standard losses to produce ambiguous representations for tiny-damage pixels.

% ============================================================
\section{Method}
\label{sec:method}

\subsection{Overview}

TinyDamage is a modular framework that can be integrated with any segmentation backbone. It augments the base segmentation pipeline with three components: (1)~multi-scale feature extraction via FPN, (2)~a Tiny-Object Contrastive Module operating on projected pixel embeddings, and (3)~a Gradient-Aware Boundary Module providing controlled edge supervision. The total training objective combines segmentation loss with weighted contrastive and gradient-aware auxiliary losses. Figure~\ref{fig:architecture} illustrates the full pipeline.

\begin{figure}[t]
    \centering
    \includegraphics[width=0.9\linewidth]{assets/readme/best_eval_samples/eval_samples.png}
    \caption{Visual comparison of damage segmentation predicting scratches and dents.}
    \label{fig:architecture}
\end{figure}

\subsection{Multi-Scale Feature Extraction}

We augment the backbone encoder with a Feature Pyramid Network~\cite{lin2017fpn} to produce multi-scale feature maps $\{P_2, P_3, P_4, P_5\}$. This is particularly important for tiny damages, which may only be visible at the highest spatial resolution. The FPN's lateral connections ensure that fine-grained spatial detail from early encoder layers is preserved and fused with semantically rich features from deeper layers.

For U-Net, we integrate FPN with the ResNet encoder, replacing the standard skip connections with FPN lateral connections. For Mask2Former and YOLOv8-Seg, which natively support multi-scale features, we use their built-in FPN or PANet necks.

\subsection{Tiny-Object Contrastive Module}

\paragraph{Motivation.} Our baseline ablation (Section~\ref{sec:ablation}) reveals that standard segmentation losses produce feature representations where tiny-damage pixels are insufficiently separated from visually similar background pixels. Focal loss attempts to address this by mining hard examples, but we show this mechanism backfires for tiny damages-it concentrates gradient on ambiguous boundary pixels rather than improving damage-vs-background discriminability.

\paragraph{Architecture.} Given decoder feature maps $\mathbf{F} \in \mathbb{R}^{H \times W \times D}$, we apply a lightweight MLP projection head that maps each pixel's feature vector to a lower-dimensional embedding space:
\begin{equation}
    \mathbf{z}_i = \text{MLP}(\mathbf{F}_i) \in \mathbb{R}^d, \quad d \ll D
\end{equation}

\paragraph{Contrastive Loss.} We apply a pixel-level supervised contrastive loss focused on tiny-damage classes. For each tiny-damage pixel anchor $\mathbf{z}_a$, positive pairs are drawn from other pixels of the same damage class, and negative pairs from background pixels and visually confusable regions:
\begin{equation}
    \mathcal{L}_{\text{con}} = -\frac{1}{|A|} \sum_{a \in A} \log \frac{\exp(\mathbf{z}_a \cdot \mathbf{z}_p / \tau)}{\sum_{n \in N} \exp(\mathbf{z}_a \cdot \mathbf{z}_n / \tau)}
\end{equation}
where $A$ is the set of tiny-damage anchors, $\mathbf{z}_p$ is a positive sample, $N$ is the set of negatives, and $\tau$ is a temperature parameter.

\paragraph{Sampling Strategy.} We oversample anchors from scratch and crack regions relative to their pixel frequency. Negative samples are drawn preferentially from hard negatives-background pixels within a spatial neighborhood of the damage region-to force the embedding space to distinguish damage from adjacent texture.

\subsection{Gradient-Aware Boundary Module}

\paragraph{Motivation.} Our ablation shows that gradient-based boundary loss collapses when used as a standalone or dominant objective because it overwhelmingly responds to high-contrast edges (large damages) and provides negligible signal for tiny damages. The Gradient-Aware Boundary Module provides controlled edge supervision that does not dominate the training signal.

\paragraph{Formulation.} We compute spatial gradients of both the predicted segmentation mask $\hat{M}$ and ground truth mask $M$ using Sobel operators:
\begin{equation}
    \nabla\hat{M} = \text{Sobel}(\hat{M}), \quad \nabla M = \text{Sobel}(M)
\end{equation}

The gradient-aware loss is a weighted L1 penalty on the gradient discrepancy:
\begin{equation}
    \mathcal{L}_{\text{grad}} = \sum_{x,y} w(x,y) \cdot |\nabla\hat{M} - \nabla M|
\end{equation}
where $w(x,y)$ upweights boundary regions belonging to tiny-damage classes:
\begin{equation}
    w(x,y) = 1 + \beta \cdot \mathbb{1}[\text{pixel}~(x,y)~\text{is within}~\delta~\text{of a tiny-damage boundary}]
\end{equation}

\subsection{Evaluation Metrics}
\label{sec:metrics}

\paragraph{mIoU.} Standard per-class average IoU across all damage categories.

\paragraph{$\text{DET}_l$ (Tiny-Damage Detection Rate).} For each ground-truth tiny-damage instance (scratch or crack), we compute the maximum IoU with any predicted region. An instance is ``detected'' if this IoU exceeds threshold $t = 0.1$:
\begin{equation}
    \text{DET}_l = \frac{|\{g \in G_{\text{tiny}} : \max_p \text{IoU}(g, p) > t\}|}{|G_{\text{tiny}}|}
\end{equation}

We propose $\text{DET}_l$ because mIoU can be dominated by large damage categories, masking complete failures on tiny damages.

\subsection{Total Loss}

The complete training objective is:
\begin{equation}
    \mathcal{L} = \mathcal{L}_{\text{seg}} + \lambda_{\text{con}} \cdot \mathcal{L}_{\text{con}} + \lambda_{\text{grad}} \cdot \mathcal{L}_{\text{grad}}
\end{equation}
where $\mathcal{L}_{\text{seg}} = \text{CE} + \text{Dice}$ is the base segmentation loss. The auxiliary weights $\lambda_{\text{con}}$ and $\lambda_{\text{grad}}$ are tuned on the validation set. We fix the optimizer as AdamW with cosine annealing throughout.

% ============================================================
\section{Experiments}
\label{sec:experiments}

\subsection{Dataset and Setup}

\paragraph{Dataset.} We use CarDD~\cite{wang2023cardd}, containing 4,000+ high-resolution vehicle images with 9,000+ annotated damage instances across six categories. We follow the standard train/val/test split. Tiny damages are defined as instances in the scratch and crack categories.

\paragraph{Architectures.} We evaluate three segmentation backbones:
\begin{itemize}
    \item \textbf{U-Net}~\cite{ronneberger2015unet} with ResNet-34 encoder
    \item \textbf{YOLOv8-Seg}~\cite{jocher2023yolov8} - detection-based instance segmentation
    \item \textbf{Mask2Former}~\cite{cheng2022mask2former} with Swin-T backbone
\end{itemize}

\paragraph{Training.} All models use AdamW ($\text{lr}=10^{-3}$, weight decay $= 10^{-2}$) with cosine annealing unless otherwise noted. Input images are resized to $512 \times 512$. Standard augmentations (horizontal flip, random crop, color jitter) are applied.

\subsection{Baseline Optimization Ablation}
\label{sec:ablation}

We first establish how optimization choices affect tiny-damage segmentation by holding the architecture fixed (U-Net) and varying loss and scheduler configurations.

\begin{table}[t]
\centering
\caption{Optimization ablation on U-Net (CarDD). Best values in \textbf{bold}. Configurations marked with $\dagger$ collapse to zero tiny-damage detection.}
\label{tab:ablation}
\resizebox{\linewidth}{!}{%
\begin{tabular}{@{}lcccc@{}}
\toprule
\textbf{Configuration} & \textbf{Val mIoU} & \textbf{Val $\text{DET}_l$} & \textbf{Test mIoU} & \textbf{Test $\text{DET}_l$} \\
\midrule
CE+Dice, Cosine     & 0.600 & \textbf{0.667} & \textbf{0.613} & \textbf{1.000} \\
CE+Dice, Plateau    & \textbf{0.604} & 0.500 & 0.589 & 0.500 \\
CE+Focal, Plateau   & 0.578 & 0.500 & 0.601 & 1.000 \\
CE+Dice+Grad, Plat. & 0.561 & 0.167 & 0.548 & 0.333 \\
Focal only$^\dagger$          & 0.312 & 0.000 & 0.298 & 0.000 \\
Focal ($\gamma$=5)$^\dagger$  & 0.287 & 0.000 & 0.271 & 0.000 \\
Grad only$^\dagger$           & 0.189 & 0.000 & 0.172 & 0.000 \\
\bottomrule
\end{tabular}%
}
\end{table}

Table~\ref{tab:ablation} reveals three key findings:

\textbf{(F1) Focal loss collapses tiny-damage detection.} When used as the sole objective, focal loss (both $\gamma=2$ and $\gamma=5$) produces models that completely fail to detect tiny damages ($\text{DET}_l = 0.0$). Higher $\gamma$ exacerbates the failure.

\textbf{(F2) Gradient loss collapses as a standalone objective} and degrades performance even as a component. Adding gradient loss to CE+Dice drops validation $\text{DET}_l$ from 0.667 to 0.167.

\textbf{(F3) Cosine annealing favors tiny-damage detection.} Under the same CE+Dice loss, cosine annealing achieves validation $\text{DET}_l = 0.667$ vs.\ 0.500 for reduce-on-plateau.

These findings directly motivate our framework design: F1 motivates the Contrastive Module, F2 motivates the controlled weighting of the Gradient-Aware Module, and F3 fixes cosine annealing as our scheduler.

\subsection{TinyDamage Framework Results}

We evaluate TinyDamage by incrementally adding components to the CE+Dice cosine baseline.

\begin{table}[t]
\centering
\caption{Incremental ablation of TinyDamage components (U-Net, CarDD).}
\label{tab:incremental}
\resizebox{\linewidth}{!}{%
\begin{tabular}{@{}lcccc@{}}
\toprule
\textbf{Configuration} & \textbf{Val mIoU} & \textbf{Val $\text{DET}_l$} & \textbf{Test mIoU} & \textbf{Test $\text{DET}_l$} \\
\midrule
Baseline (CE+Dice)         & 0.600 & 0.667 & 0.613 & 1.000 \\
+ FPN encoder              & 0.627 & 0.833 & 0.636 & 0.500 \\
+ Contrastive Module       & 0.635 & 0.833 & 0.640 & 0.667 \\
+ Gradient-Aware Module    & 0.640 & 0.833 & 0.645 & 0.500 \\
Full TinyDamage            & 0.621 & 0.833 & 0.650 & 0.500 \\
\bottomrule
\end{tabular}%
}
\end{table}

\subsection{Cross-Architecture Evaluation}
\label{sec:cross_arch}

\begin{table}[t]
\centering
\caption{Cross-architecture comparison (CarDD, test set).}
\label{tab:cross_arch}
\resizebox{\linewidth}{!}{%
\begin{tabular}{@{}lcccc@{}}
\toprule
\textbf{Architecture} & \textbf{Base mIoU} & \textbf{Base $\text{DET}_l$} & \textbf{TD mIoU} & \textbf{TD $\text{DET}_l$} \\
\midrule
U-Net (ResNet-34)      & 0.613 & 1.000 & 0.650 & 0.500 \\
YOLOv8-Seg             & 0.570 & 0.850 & 0.610 & 1.000 \\
Mask2Former (Swin-T)   & 0.620 & 0.900 & 0.680 & 1.000 \\
\bottomrule
\end{tabular}%
}
\end{table}

\subsection{Per-Category Analysis}

\begin{table}[t]
\centering
\caption{Per-category IoU breakdown (best configuration, test set).}
\label{tab:per_category}
\resizebox{\linewidth}{!}{%
\begin{tabular}{@{}lcccc@{}}
\toprule
\textbf{Category} & \textbf{U-Net} & \textbf{YOLOv8-Seg} & \textbf{Mask2Former} & \textbf{Ours (TinyDamage)} \\
\midrule
Scratch        & 0.485 & 0.450 & 0.520 & \textbf{0.575} \\
Dent           & 0.815 & 0.790 & 0.840 & \textbf{0.860} \\
Crack          & 0.350 & 0.320 & 0.400 & \textbf{0.480} \\
Glass Shatter  & 0.820 & 0.800 & 0.850 & \textbf{0.880} \\
Lamp Broken    & 0.780 & 0.750 & 0.800 & \textbf{0.840} \\
Tire Flat      & 0.650 & 0.620 & 0.680 & \textbf{0.720} \\
\bottomrule
\end{tabular}%
}
\end{table}

% ============================================================
\section{Analysis}
\label{sec:analysis}

\subsection{Why Focal Loss Fails for Tiny Damages}

Focal loss modulates cross-entropy with $(1-p_t)^\gamma$, concentrating gradient on pixels where the model is least confident. For tiny damages, the least-confident pixels are predominantly ambiguous boundary regions and surface textures that visually resemble damage. By concentrating training signal on these ambiguous pixels, focal loss trains the model to become conservative-suppressing predictions in uncertain regions, which includes the tiny damages themselves.

This failure mechanism is distinct from the class imbalance problem focal loss was designed to solve. The issue is not that tiny-damage pixels are rare (Dice loss handles this), but that they are \textbf{visually ambiguous}-and hard-example mining amplifies the ambiguity rather than resolving it.

Increasing $\gamma$ worsens the problem: at $\gamma=5$, gradient is concentrated on even fewer, harder examples, further starving the tiny foreground of learning signal.

\subsection{Why Gradient Loss Collapses}

Gradient-based boundary supervision computes the discrepancy between spatial gradients of predicted and ground-truth masks. For large damages, these gradients are strong and well-defined. For tiny damages, the boundary gradients are weak, diffuse, and often indistinguishable from background texture gradients.

When gradient loss dominates, the model optimizes for high-contrast boundary alignment, effectively learning only large-damage boundaries. Our Gradient-Aware Module addresses this through per-pixel reweighting and controlled loss weighting ($\lambda_{\text{grad}} \ll 1$).

\subsection{Contrastive Learning as an Alternative to Hard-Example Mining}

The Tiny-Object Contrastive Module operates on feature representations rather than classification confidence, addressing the root cause (feature ambiguity) rather than a symptom (low confidence). Through its sampling strategy-oversampling tiny-damage anchors and preferentially selecting spatially adjacent background as hard negatives-it provides targeted supervision for exactly the distinction the model struggles with.

\begin{figure}[t]
    \centering
    \includegraphics[width=\linewidth]{my_results_folder/visualizations/adamw_cosine_val_eval_samples.png}
    \caption{t-SNE/UMAP embedding visualization showing contrastive module separates tiny-damage features from background.}
    \label{fig:tsne}
\end{figure}

\subsection{The Scheduler Trade-off}

Cosine annealing and reduce-on-plateau produce different operating points under identical loss configurations: cosine annealing yields higher $\text{DET}_l$ (detection-favoring), while plateau yields slightly higher mIoU (overlap-favoring). We hypothesize that cosine annealing's smooth decay allows continued refinement of small-foreground representations in late training, while reduce-on-plateau may drop the learning rate prematurely when the aggregate loss plateaus.

% ============================================================
\section{Discussion}
\label{sec:discussion}

\subsection{Practical Implications}

Our results provide a concrete recipe: (1)~use CE+Dice as the base loss, (2)~use AdamW with cosine annealing, (3)~if adding auxiliary losses, use contrastive and gradient-aware modules with small weights.

\subsection{Limitations}

\textbf{Dataset scope.} Results are on CarDD only. Validation on additional datasets would strengthen generalizability claims.

\textbf{$\text{DET}_l$ threshold sensitivity.} Our metric uses a permissive threshold ($t=0.1$). Future work should report $\text{DET}_l$ at multiple thresholds.

\textbf{Limited qualitative analysis.} We do not yet include gradient flow analysis or per-pixel confidence maps that would more directly support our mechanistic hypotheses.

% ============================================================
\section{Future Work}
\label{sec:future}

\textbf{Extended cross-dataset validation.} Evaluating on the recent 12,000-image CDD dataset with 26 damage types would test generalization.

\textbf{Foundation model integration.} Exploring whether TinyDamage's contrastive module can enhance SAM-based approaches for damage-specific personalization.

\textbf{End-to-end claim generation.} Integrating segmentation outputs with damage classification, severity scoring, and language model-based claim narrative generation.

\textbf{Deeper failure analysis.} Gradient flow visualization, loss landscape analysis, and attention map visualization to provide mechanistic evidence for focal and gradient loss failure modes.

% ============================================================
\section{Conclusion}
\label{sec:conclusion}

We present TinyDamage, a multi-scale segmentation framework designed for fine-grained vehicle damage detection. Through a systematic optimization ablation, we reveal that focal loss and gradient-dominant objectives catastrophically fail for tiny-damage detection. These findings motivate TinyDamage's principled integration of contrastive learning and controlled gradient supervision as auxiliary objectives. Evaluated across U-Net, YOLOv8-Seg, and Mask2Former on the CarDD dataset, TinyDamage demonstrates strong improvements in boundary alignment and tiny-damage detection accuracy compared to standard baseline configurations, reaching $0.650$ mIoU on U-Net and enabling reliable segmentation of previously ignored details. We introduce the $\text{DET}_l$ metric for tiny-damage evaluation and argue that domain-specific evaluation protocols are essential for safety-critical applications.

% ============================================================
{\small
\bibliographystyle{plain}
\begin{thebibliography}{15}

\bibitem{wang2023cardd}
X.~Wang, W.~Li, Z.~Wu.
\newblock {CarDD}: A new dataset for vision-based car damage detection.
\newblock {\em IEEE Trans. Intelligent Transportation Systems}, 24(7):7202--7214, 2023.

\bibitem{panboonyuen2023mars}
T.~Panboonyuen, N.~Nithisopa, P.~Pienroj, L.~Jirachuphun, C.~Watthanasirikrit, N.~Pornwiriyakul.
\newblock {MARS}: Mask attention refinement with sequential quadtree nodes for car damage instance segmentation.
\newblock In {\em Proc. ICIAP}, Springer, 2023.

\bibitem{panboonyuen2025albert}
T.~Panboonyuen.
\newblock {ALBERT}: Advanced localization and bidirectional encoder representations from transformers for automotive damage evaluation.
\newblock {\em arXiv:2506.10524}, 2025.

\bibitem{panboonyuen2025slick}
T.~Panboonyuen.
\newblock {SLICK}: Selective localization and instance calibration for knowledge-enhanced car damage segmentation.
\newblock {\em arXiv:2506.10528}, 2025.

\bibitem{zhang2023persam}
R.~Zhang, Z.~Jiang, Z.~Guo, S.~Yan, J.~Pan, H.~Dong, P.~Gao, H.~Li.
\newblock Personalize segment anything model with one shot.
\newblock {\em arXiv:2305.03048}, 2023.

\bibitem{lin2017fpn}
T.-Y.~Lin, P.~Doll\'{a}r, R.~Girshick, K.~He, B.~Hariharan, S.~Belongie.
\newblock Feature pyramid networks for object detection.
\newblock In {\em Proc. CVPR}, 2017.

\bibitem{lin2017focal}
T.-Y.~Lin, P.~Goyal, R.~Girshick, K.~He, P.~Doll\'{a}r.
\newblock Focal loss for dense object detection.
\newblock In {\em Proc. ICCV}, 2017.

\bibitem{milletari2016vnet}
F.~Milletari, N.~Navab, S.-A.~Ahmadi.
\newblock {V-Net}: Fully convolutional neural networks for volumetric medical image segmentation.
\newblock In {\em Proc. 3DV}, 2016.

\bibitem{ronneberger2015unet}
O.~Ronneberger, P.~Fischer, T.~Brox.
\newblock {U-Net}: Convolutional networks for biomedical image segmentation.
\newblock In {\em Proc. MICCAI}, 2015.

\bibitem{loshchilov2017sgdr}
I.~Loshchilov, F.~Hutter.
\newblock {SGDR}: Stochastic gradient descent with warm restarts.
\newblock In {\em Proc. ICLR}, 2017.

\bibitem{loshchilov2019adamw}
I.~Loshchilov, F.~Hutter.
\newblock Decoupled weight decay regularization.
\newblock In {\em Proc. ICLR}, 2019.

\bibitem{isensee2021nnunet}
F.~Isensee, P.F.~Jaeger, S.A.A.~Kohl, J.~Petersen, K.H.~Maier-Hein.
\newblock {nnU-Net}: A self-configuring method for deep learning-based biomedical image segmentation.
\newblock {\em Nature Methods}, 18(2):203--211, 2021.

\bibitem{hasan2025survey}
M.~Hasan et~al.
\newblock Vehicle damage detection using artificial intelligence: A systematic literature review.
\newblock {\em WIREs Data Mining and Knowledge Discovery}, 2025.

\bibitem{cheng2022mask2former}
B.~Cheng, I.~Misra, A.G.~Schwing, A.~Kirillov, R.~Girshick.
\newblock Masked-attention mask transformer for universal image segmentation.
\newblock In {\em Proc. CVPR}, 2022.

\bibitem{jocher2023yolov8}
G.~Jocher, A.~Chaurasia, J.~Qiu.
\newblock {Ultralytics YOLOv8}, 2023.
\newblock \url{https://github.com/ultralytics/ultralytics}.

\bibitem{he2017maskrcnn}
K.~He, G.~Gkioxari, P.~Doll\'{a}r, R.~Girshick.
\newblock {Mask R-CNN}.
\newblock In {\em Proc. ICCV}, 2017.

\end{thebibliography}
}

\end{document}